\begin{problem}[Matrix Inverses and Matrix Equations]
Let $A$ and $B$ be $2 \times 2$ matrices defined as $A = \begin{bmatrix} k & 2 \\ -1 & 1 \end{bmatrix}$ and $B = \begin{bmatrix} 1 & -1 \\ 2 & 3 \end{bmatrix}$, where $k$ is a real constant.

\begin{enumerate}[label=\textbf{(\alph*)}]
    \item Determine the determinant of $A$ in terms of $k$, and hence state the value of $k$ for which matrix $A$ is singular.
    \item Let $k = 3$. Determine the inverse matrices $A^{-1}$ and $B^{-1}$.
    \item For $k = 3$, calculate the matrix product $AB$ and hence determine $(AB)^{-1}$.
    \item Evaluate the product $B^{-1}A^{-1}$ to verify the property $(AB)^{-1} = B^{-1}A^{-1}$. Hence, use this property to make the matrix $X$ the subject of the matrix equation $ABX = C$, where $C$ is a $2 \times 2$ matrix.
\end{enumerate}
\end{problem}

\begin{hint}
\begin{itemize}
    \item \textbf{Part (a):} The determinant of a $2 \times 2$ matrix $\begin{bmatrix} a & b \\ c & d \end{bmatrix}$ is $ad - bc$. A matrix is singular (has no inverse) when its determinant equals zero.
    \item \textbf{Part (b):} Use the formula for the inverse of a $2 \times 2$ matrix: $\frac{1}{\det(M)} \begin{bmatrix} d & -b \\ -c & a \end{bmatrix}$.
    \item \textbf{Part (c):} Multiply matrix $A$ by matrix $B$ using the standard row-by-column method. Then find the inverse of the resulting $2 \times 2$ matrix.
    \item \textbf{Part (d):} Multiply $B^{-1}$ (on the left) by $A^{-1}$ (on the right) and factor out the scalar multipliers. To solve $ABX = C$, remember that you must multiply both sides of the equation by the inverse of $AB$ on the left.
\end{itemize}
\end{hint}

\begin{solution}
\textbf{(a)} \\
$\det(A) = (k)(1) - (2)(-1) = k + 2$. \\
For matrix $A$ to be singular, $\det(A) = 0$. Therefore, $k = -2$.

\medskip
\textbf{(b)} \\
Given $k = 3$, $A = \begin{bmatrix} 3 & 2 \\ -1 & 1 \end{bmatrix}$. \\
$\det(A) = 3 + 2 = 5$. Thus, $A^{-1} = \frac{1}{5}\begin{bmatrix} 1 & -2 \\ 1 & 3 \end{bmatrix}$. \\
$\det(B) = (1)(3) - (-1)(2) = 3 + 2 = 5$. Thus, $B^{-1} = \frac{1}{5}\begin{bmatrix} 3 & 1 \\ -2 & 1 \end{bmatrix}$.

\medskip
\textbf{(c)} \\
$AB = \begin{bmatrix} 3 & 2 \\ -1 & 1 \end{bmatrix} \begin{bmatrix} 1 & -1 \\ 2 & 3 \end{bmatrix} = \begin{bmatrix} 3(1)+2(2) & 3(-1)+2(3) \\ -1(1)+1(2) & -1(-1)+1(3) \end{bmatrix} = \begin{bmatrix} 7 & 3 \\ 1 & 4 \end{bmatrix}$. \\
$\det(AB) = (7)(4) - (3)(1) = 28 - 3 = 25$. \\
$(AB)^{-1} = \frac{1}{25}\begin{bmatrix} 4 & -3 \\ -1 & 7 \end{bmatrix}$.

\medskip
\textbf{(d)} \\
$B^{-1}A^{-1} = \left( \frac{1}{5}\begin{bmatrix} 3 & 1 \\ -2 & 1 \end{bmatrix} \right) \left( \frac{1}{5}\begin{bmatrix} 1 & -2 \\ 1 & 3 \end{bmatrix} \right) = \frac{1}{25} \begin{bmatrix} 3(1)+1(1) & 3(-2)+1(3) \\ -2(1)+1(1) & -2(-2)+1(3) \end{bmatrix} = \frac{1}{25}\begin{bmatrix} 4 & -3 \\ -1 & 7 \end{bmatrix}$. \\
This verifies that $(AB)^{-1} = B^{-1}A^{-1}$. \\
To solve $ABX = C$: \\
$(AB)^{-1}(AB)X = (AB)^{-1}C$ \\
$IX = (AB)^{-1}C$ \\
$X = B^{-1}A^{-1}C$.
\end{solution}

\begin{takeaways}
\item \textbf{The "Socks and Shoes" Property:} The inverse of a matrix product reverses the order of the original matrices: $(AB)^{-1} = B^{-1}A^{-1}$.
\item \textbf{Matrix Algebra in Equations:} When isolating a matrix variable (like $X$ in $ABX=C$), you must respect the non-commutative nature of matrices by multiplying by the inverse on the correct side (in this case, pre-multiplying on the left).
\item \textbf{Determinant of a Product:} Notice that $\det(AB) = 25$, while $\det(A) = 5$ and $\det(B) = 5$. This verifies another key property: $\det(AB) = \det(A) \times \det(B)$.
\end{takeaways}


\begin{problem}[Matrix Inverses and the Zero Matrix]
Let $A, B, C,$ and $D$ be $2 \times 2$ matrices, where $O$ is the zero matrix and $I$ is the identity matrix.

\begin{enumerate}[label=\textbf{(\alph*)}]
    \item Assume $A$ is an invertible matrix and that $AB = O$. Prove algebraically that $B$ must be the zero matrix.
    \item Suppose instead that $C$ and $D$ are both strictly non-zero matrices, yet their product is the zero matrix ($CD = O$). By referencing your proof in part (a), deduce whether $C$ can be an invertible matrix. Support your answer with logical reasoning.
    \item A matrix $M$ is defined as being self-inverse if $M^{-1} = M$. By showing this implies $M^2 = I$, determine the algebraic condition on real numbers $a, b,$ and $c$ such that $M = \begin{bmatrix} a & b \\ c & -a \end{bmatrix}$ is a self-inverse matrix.
    \item Using your condition from part (c), construct a specific matrix $M$ where $a = 3$, and $b$ and $c$ are integers. Verify that your chosen matrix is self-inverse by explicitly evaluating $M^2$.
\end{enumerate}
\end{problem}

\begin{hint}
\begin{itemize}
    \item \textbf{Part (a):} Start with the given equation $AB = O$. Since $A$ is invertible, you are permitted to pre-multiply both sides of the equation by $A^{-1}$. What does $A^{-1}A$ simplify to?
    \item \textbf{Part (b):} Use a proof by contradiction. What would happen if $C$ \textit{was} invertible? Apply the exact same logic established in part (a) and compare the result to the condition that $D$ is a non-zero matrix.
    \item \textbf{Part (c):} Multiply both sides of $M^{-1} = M$ by $M$. Then, square the given matrix $M = \begin{bmatrix} a & b \\ c & -a \end{bmatrix}$ using standard row-by-column multiplication and equate the resulting elements to the corresponding elements of the identity matrix, $I$.
    \item \textbf{Part (d):} Substitute $a=3$ into your condition from part (c). Find two integers $b$ and $c$ that multiply together to satisfy the equation, construct the matrix, and manually multiply $M \times M$ to check it yields $I$.
\end{itemize}
\end{hint}

\begin{solution}
\textbf{(a)} \\
Given $AB = O$ and $A$ is invertible, $A^{-1}$ exists. Pre-multiplying both sides by $A^{-1}$: \\
$A^{-1}(AB) = A^{-1}O$ \\
$(A^{-1}A)B = O$ \\
$IB = O \implies B = O$.

\medskip
\textbf{(b)} \\
Assume for the sake of contradiction that $C$ is invertible. \\
Then, given $CD = O$, we could pre-multiply by $C^{-1}$ to get $C^{-1}(CD) = C^{-1}O \implies D = O$. \\
However, the problem states that $D$ is a non-zero matrix. This is a contradiction. Therefore, the assumption that $C$ is invertible must be false; $C$ must be a singular matrix.

\medskip
\textbf{(c)} \\
If $M^{-1} = M$, then multiplying both sides by $M$ gives $MM^{-1} = MM \implies I = M^2$. \\
Evaluating $M^2$: \\
$M^2 = \begin{bmatrix} a & b \\ c & -a \end{bmatrix} \begin{bmatrix} a & b \\ c & -a \end{bmatrix} = \begin{bmatrix} a^2+bc & ab-ab \\ ac-ac & bc+(-a)^2 \end{bmatrix} = \begin{bmatrix} a^2+bc & 0 \\ 0 & a^2+bc \end{bmatrix}$. \\
For $M^2 = I = \begin{bmatrix} 1 & 0 \\ 0 & 1 \end{bmatrix}$, we must have the condition: $a^2 + bc = 1$.

\medskip
\textbf{(d)} \\
Let $a=3$. From part (c), $3^2 + bc = 1 \implies 9 + bc = 1 \implies bc = -8$. \\
Choose integers $b = 2$ and $c = -4$ (other integer pairs like $4, -2$ are also valid). \\
$M = \begin{bmatrix} 3 & 2 \\ -4 & -3 \end{bmatrix}$. \\
Verification: \\
$M^2 = \begin{bmatrix} 3 & 2 \\ -4 & -3 \end{bmatrix} \begin{bmatrix} 3 & 2 \\ -4 & -3 \end{bmatrix} = \begin{bmatrix} 9-8 & 6-6 \\ -12+12 & -8+9 \end{bmatrix} = \begin{bmatrix} 1 & 0 \\ 0 & 1 \end{bmatrix} = I$.
\end{solution}

\begin{takeaways}
\item \textbf{Zero Product Property Failure:} In real number algebra, $xy=0$ implies $x=0$ or $y=0$. In matrix algebra, $CD=O$ does \textit{not} necessarily mean $C=O$ or $D=O$. Non-zero matrices can multiply to give the zero matrix.
\item \textbf{Singularity is Key:} For two non-zero matrices to multiply to the zero matrix, they must both be singular (determinant of zero). If either were invertible, the other would be forced to be the zero matrix.
\item \textbf{Self-Inverse Matrices:} Matrices where $M^{-1} = M$ exist in infinitely many non-trivial forms beyond just the identity matrix. Geometrically, self-inverse (or involutory) matrices often represent reflections, where applying the transformation twice returns all coordinates to their original positions.
\end{takeaways}


\begin{problem}[Matrix Methods for Simultaneous Equations]
Consider the following system of simultaneous linear equations with real parameter $k$:
\begin{align*}
(k - 1)x + 3y &= k + 2 \\
2x + (k - 2)y &= 4
\end{align*}

\begin{enumerate}[label=\textbf{(\alph*)}]
    \item Express this system as a matrix equation of the form $AX = B$, and find a simplified quadratic expression for the determinant of the coefficient matrix $A$.
    \item Determine the values of $k$ for which the matrix $A$ is singular (i.e., does not have an inverse).
    \item For each value of $k$ found in part (b), determine whether the system has no solutions or infinitely many solutions. Provide a brief geometric justification for your answers.
    \item Assuming the matrix $A$ is invertible, use the matrix inverse method to find the unique solution for $x$ and $y$ in terms of $k$, expressing your final answers in their simplest rational form.
\end{enumerate}
\end{problem}

\begin{hint}
\begin{itemize}
    \item \textbf{Part (a):} The coefficient matrix $A$ contains the multipliers for $x$ and $y$. The determinant of a $2 \times 2$ matrix $\begin{bmatrix} a & b \\ c & d \end{bmatrix}$ is calculated as $ad - bc$.
    \item \textbf{Part (b):} A matrix is singular when its determinant is exactly zero. Set your expression from part (a) to zero and solve the resulting quadratic equation for $k$.
    \item \textbf{Part (c):} Substitute your $k$ values back into the original system of equations. If the two equations represent distinct parallel lines (e.g., yielding a false statement like $0=5$), there is no solution. If they represent the exact same line, there are infinitely many solutions.
    \item \textbf{Part (d):} Use the inverse matrix formula $X = A^{-1}B = \frac{1}{\det(A)} \begin{bmatrix} d & -b \\ -c & a \end{bmatrix} B$. After multiplying the matrices, factor the resulting numerators to see if they cancel with the determinant in the denominator.
\end{itemize}
\end{hint}

\begin{solution}
\textbf{(a)} \\
Matrix equation: 
$\begin{bmatrix} k-1 & 3 \\ 2 & k-2 \end{bmatrix} \begin{bmatrix} x \\ y \end{bmatrix} = \begin{bmatrix} k+2 \\ 4 \end{bmatrix}$ \\
$\det(A) = (k-1)(k-2) - (3)(2) = k^2 - 3k + 2 - 6 = k^2 - 3k - 4$. \\
Factored form: $\det(A) = (k-4)(k+1)$.

\medskip
\textbf{(b)} \\
Matrix $A$ is singular when $\det(A) = 0$. \\
$k^2 - 3k - 4 = 0 \implies (k-4)(k+1) = 0$. \\
Therefore, $A$ is singular when $k = 4$ or $k = -1$.

\medskip
\textbf{(c)} \\
\textbf{Case $k = 4$:} \\
Equations become $3x + 3y = 6$ (which simplifies to $x+y=2$) and $2x + 2y = 4$ (which simplifies to $x+y=2$). \\
Because the equations represent coincident (identical) lines, there are \textbf{infinitely many solutions}. \\
\textbf{Case $k = -1$:} \\
Equations become $-2x + 3y = 1$ and $2x - 3y = 4$. \\
Adding the two equations yields $0 = 5$, which is mathematically inconsistent. Because the equations represent distinct parallel lines, there are \textbf{no solutions}.

\medskip
\textbf{(d)} \\
For $k \neq 4$ and $k \neq -1$, $A$ is invertible. \\
$A^{-1} = \frac{1}{(k-4)(k+1)} \begin{bmatrix} k-2 & -3 \\ -2 & k-1 \end{bmatrix}$ \\
$X = \frac{1}{(k-4)(k+1)} \begin{bmatrix} k-2 & -3 \\ -2 & k-1 \end{bmatrix} \begin{bmatrix} k+2 \\ 4 \end{bmatrix}$ \\
$x = \frac{(k-2)(k+2) - 12}{(k-4)(k+1)} = \frac{k^2 - 4 - 12}{(k-4)(k+1)} = \frac{k^2 - 16}{(k-4)(k+1)} = \frac{(k-4)(k+4)}{(k-4)(k+1)} = \frac{k+4}{k+1}$ \\
$y = \frac{-2(k+2) + 4(k-1)}{(k-4)(k+1)} = \frac{-2k - 4 + 4k - 4}{(k-4)(k+1)} = \frac{2k - 8}{(k-4)(k+1)} = \frac{2(k-4)}{(k-4)(k+1)} = \frac{2}{k+1}$ \\
Unique solution: $x = \frac{k+4}{k+1}$ and $y = \frac{2}{k+1}$.
\end{solution}

\begin{takeaways}
\item \textbf{Determinant as a Geometric Indicator:} Just as the discriminant reveals the nature of roots in a quadratic, the determinant of a coefficient matrix dictates the geometry of linear systems. A non-zero determinant guarantees intersecting lines (one unique solution).
\item \textbf{Interpreting Singularity:} A singular matrix ($\det = 0$) means the matrix transformation collapses 2D space. For a system of equations, this means the lines share the exact same gradient, resulting in either parallel lines (no intersection) or coincident lines (infinite intersections).
\item \textbf{Algebraic Simplification:} When solving parametrically with the matrix inverse method, finding the solution often involves polynomial division or factoring. The non-zero roots of the determinant will frequently appear as cancellable factors in the numerators of your $x$ and $y$ expressions.
\end{takeaways}
